% Build the layer, Chapter 0: two increments, and a check where the number
% changes. Inserted at the end of \S\ref{pre:sec:kernel-gap}, after the
% subsection on why a model cannot be its own kernel.
%
% No code or output in this file is typed. \pdexcerpt and \pdrun set files
% under excerpts/, which the kernel's tools/excerpts.mjs cuts from its
% source and its recorded runs (see figures/pd-listing.tex). The figure
% plots four of those recorded runs as they were printed. To change a
% listing or a curve, change the kernel and cut again. The prose names
% functions, never line numbers, so that a listing which moves in its file
% does not make a sentence wrong.
\subsection{Build the layer: two increments, both kept}
\label{sec:prereq-build}

\noindent\emph{Express lane: the chapter's opening example is built and run.
A worker that checks its own reading still loses an update. A check made
where the number changes loses none, in any order the workers can move in.
A simulation then follows each design through 600 rounds.}

The chapter opened with two workers and a counter, and one update went
missing. Theorem~\ref{thm:kernel-gap} has since said where the check must
sit, and \S\ref{sec:prereq-no-self-kernel} has argued that the worker cannot
be the one to make it. This section runs both claims. The code is part of
\pdkernelname{}, a small kernel written for this book, which gains a layer
in each chapter. Its source is not yet published. The layer here needs none
of the kernel's machinery. It is a number held in memory, and a rule that
decides whether a worker's write to it lands.

Every chapter after this one has a section of the same shape, in four beats.
\emph{Naive} builds the design most people reach for first, and builds it
with care. \emph{Break} is a run that makes it fail, the same way every
time. \emph{Fix} is the chapter's mechanism. \emph{Verify} runs the break
against the fix, and then against the fix with its guard taken out, to show
that the guard is what stopped it.

None of the code below was typed into the book. A script cuts each listing
from the kernel's source. The listing's head names the file and the lines,
and the numbers beside the code are the file's own. Each run is a recording,
and one of the kernel's tests compares a fresh run with it. The kernel's
repository also checks these listings against its source and fails when the
two differ.

\subsubsection{Naive: read, then write}

The number is kept by a small object that we will call the \emph{cell}. It
holds the number and a \emph{version}, a count of the writes that have
landed. A worker can \texttt{read} the cell, which returns both, and can ask
it to \texttt{put} a new number. The cell takes one request at a time and
finishes it before it looks at the next, so two writes never land at once.
That is the queue of the opening example.

A worker that adds one does the obvious thing
(Listing~\ref{alg:prereq-build-blind}).

\pdexcerpt[caption={The obvious worker. It reads the number and, after
thinking, writes the number plus one.},label={alg:prereq-build-blind}]{ch00-blind}

The line \texttt{yield* thinking(think)} is why the worker is written as a
generator. Everything a worker does between two \texttt{yield}s is one
\emph{turn}: it runs to its end with nothing else running. At a
\texttt{yield} the worker stops, and another worker may take a turn before
it goes on. With \texttt{think} at 1, as it is until the long run below,
the blind worker takes 2 turns: it reads in the first and writes in the
second. A real worker has no such marks in its source, and it stops all
the same, whenever it waits for a model to answer or for a request to come
back. Writing the stops down lets us choose who moves next. That choice is a
function, the \pdterm{scheduler}, and the loop that applies it is short
(Listing~\ref{alg:prereq-build-run}).

\pdexcerpt[caption={One turn at a time. The scheduler \texttt{next} picks
who moves.},label={alg:prereq-build-run}]{ch00-run}

\subsubsection{Break: the workers take turns}

Start the number at 0 and let Heron and Wren each add one, moving in turn
with Heron first.

\pdrun[caption={The break. Both workers read 0, and both write
1.},label={alg:prereq-build-run-blind}]{increments-blind}

Both writes were taken. The version went to 2, so two writes landed, one
after the other, and the number is 1. No two writes collided, and each
worker wrote exactly what it had worked out. Wren worked it out from a
reading that Heron's write had already made old.

One bad order could be bad luck. Two workers with two turns each can move in
only six orders, so \pdkernelname{} tries them all
(Listing~\ref{alg:prereq-build-orders-blind}).

\pdrun[caption={Every order of the blind workers' turns: who moved, turn by
turn, and the number at the end.},label={alg:prereq-build-orders-blind}]{orders-blind}

Four of the six lose an update. The two that keep both are the orders in
which one worker finishes before the other starts. This is part~(2) of
Theorem~\ref{thm:kernel-gap}. The queue puts the writes in order and says
nothing about the readings they were worked out from.

\subsubsection{Break again: the worker checks its own reading}

The first repair most people try is to have the worker look again. Just
before it writes, it reads the cell a second time. If the version has moved,
it throws its work away and starts over
(Listing~\ref{alg:prereq-build-look-twice}).

\pdexcerpt[caption={Looking twice. The worker compares a second reading
with its first, and writes if they agree.},label={alg:prereq-build-look-twice}]{ch00-look-twice}

\pdrun[caption={Looking twice, with the workers taking turns. Steps 3 and 4
are the second looks.},label={alg:prereq-build-run-look-twice}]{increments-look-twice}

Each worker looked again, saw version 0 as before, and wrote. The second
look caught nothing because, when it was made, there was nothing to catch:
the other write came after it. The \texttt{yield} between the look and the
\texttt{put} is not a slip that better code would remove. The look and the
write are two requests, and the cell is free to serve another worker between
them. A worker can move its check closer to its write. It cannot move the
check into the write, because the write happens in the cell.

This is the first argument of \S\ref{sec:prereq-no-self-kernel} as a
program. A model that reviews its own tool call is looking twice: it judges
the state that its prompt described, and the call takes effect later.
\pdchaplink{swk}{Chapter~1} meets the same gap with a lease and a clock in
it (\weld{swk:sec:swk-build}{lease}). \emph{Now you try:}
Exercise~\ref{ex:prereq-build-orders} asks which orders defeat the second
look.

\subsubsection{Fix: the write says what it read}

The fix changes one call. A worker's \texttt{put} now carries the version
that the worker read, and a refused \texttt{put} sends the worker back to
read again (Listing~\ref{alg:prereq-build-checked}). This worker is not more
careful than the last one. It does less. What changed is on the other side
of the call. The cell hands each \texttt{put} to the kernel's
\pdterm{decision core}, the one function through which every proposal in
this book's code is decided. The cell passes its own state at that moment,
and keeps what the core returns
(Listing~\ref{alg:prereq-build-shared-put}).

\Needspace{12\baselineskip}
\pdexcerpt[caption={The checked worker. Its write names the version it
read.},label={alg:prereq-build-checked}]{ch00-checked}

\Needspace{14\baselineskip}
\pdexcerpt[caption={The cell's \texttt{put}. It asks the core, and stores
the answer before it returns.},label={alg:prereq-build-shared-put}]{ch00-shared-put}

The decision for a \texttt{put} is a few lines of Rust
(Listing~\ref{alg:prereq-build-put}). Unless the cell's policy is the blind
one of the first two designs, it calls the guard, and then it counts the
write.

\Needspace{16\baselineskip}
\pdexcerpt[caption={The decision for a write to the
cell.},label={alg:prereq-build-put}]{ch00-put}

The guard is \texttt{current} (Listing~\ref{alg:prereq-build-current}). It
compares the version that the write names with the version the cell is at.
If they differ, something was written after the worker read, and the write
is refused with both numbers in the reason. A write that names no version is
refused as well. The cell cannot tell an old reading from no reading, and a
check that a caller can skip by leaving out a field does not cover every
write.

\Needspace{14\baselineskip}
\pdexcerpt[caption={The guard: nothing has been written since the writer
read.},label={alg:prereq-build-current}]{ch00-current}

The comparison and the change are one step. \texttt{decide} returns the new
cell or a reason, and the cell's \texttt{put} stores the answer before it
returns, so no worker's turn can fall between the two. Part~(1) of
Theorem~\ref{thm:kernel-gap} asks for exactly this boundary.

\keyidea[A write is good for the version it read]{A worker cannot know,
when it writes, whether what it read is still current. The cell can. So the
write names the version it was worked out from, and the cell compares that
with its own in the same step as the change. The fencing epoch of Chapter~1
is this idea applied to who holds a file
(\weld{swk:sec:swk-build}{fence}), and its rule for a test run applies it
to a result (\weld{swk:sec:swk-roads}{revision}).}

The idea is old. Databases have called it optimistic concurrency control
since Kung and Robinson~\pdcite{kung1981optimistic}: let the work go ahead
without a lock, and validate it at the moment of writing. HTTP carries it as
the \texttt{If-Match} header, where the version is an entity tag and the
refusal is status 412~\pdcite{rfc9110}.

\subsubsection{Verify: every order, then the guard removed}

The same two workers, taking turns as before:

\Needspace{16\baselineskip}
\pdrun[caption={The fix. Wren's write from version 0 is refused, and Wren
starts over.},label={alg:prereq-build-run-checked}]{increments-checked}

Step 4 is the refusal. Wren read at version 0, and by the time its write
arrived the cell was at version 1. Wren reads again and writes 2 from the 1
it finds.

One good run proves as little as one bad run did, so here is every order
for each design (Listing~\ref{alg:prereq-build-orders}).

\pdrun[caption={Every order, for each design, with two
workers.},label={alg:prereq-build-orders}]{orders}

The checked design loses no update in any order. Looking twice has 20 orders
because each worker now takes three turns, and it loses an update in 12 of
them. The last column weighs the orders by how likely they are when each
turn goes to either worker with equal chance. Under that rule some orders
are likelier than others, which is why 12 orders of 20 come to a chance of
3/8.

A third worker changes the counts and leaves the conclusion where it was
(Listing~\ref{alg:prereq-build-orders-three}). Looking twice gets worse with
more workers, and the checked design still loses nothing in any of its 198
orders.

\pdrunlines[caption={Every order with three
workers.},label={alg:prereq-build-orders-three}]{orders-three}{1-8}{1\textendash 8}

Last, the guard. \pdkernelname{} keeps a second build of its core in which
\texttt{current} returns at once, at the line marked \texttt{MUTANT} in
Listing~\ref{alg:prereq-build-current}. With that build in place of the real
one, the checked design loses the same four orders that the blind design
loses (Listing~\ref{alg:prereq-build-orders-mutant}).

\pdrunlines[caption={The fix with its guard
removed.},label={alg:prereq-build-orders-mutant}]{orders-blind-put}{1-9}{1\textendash 9}

So the comparison in \texttt{current} is what keeps the update, and nothing
else about the checked design does.

\subsubsection{The long run: how long a job stays safe}

A chance of 3/8 is smaller than a chance of 1/2.
Theorem~\ref{thm:non-determinism} says what a smaller chance buys over a
long job. If every step has at least a chance $\epsilon$ of letting a bad
effect through, the chance of getting through $H$ steps is at most
$(1-\epsilon)^H$, and that goes to zero for any $\epsilon$ above zero. We
can watch it happen.

\paragraph{The workload.}
This part is a model, and its numbers are our choice. A job is 600 rounds.
In each round Heron and Wren each add one to the number. In most rounds one
of them finishes before the other starts, and nothing can be lost. In 1
round of 50 their turns overlap, and each turn then goes to either worker
with equal chance. Listing~\ref{alg:prereq-build-job} is one job. After
every round it compares the number with the count of increments so far, and
it returns the first round in which the number has fallen behind.

\pdexcerpt[caption={One job. In a round that overlaps, the scheduler is
\texttt{atRandom}.},label={alg:prereq-build-job}]{ch00-job}

A worker thinks for 8 turns between reading and writing, because working
out an edit takes far longer than looking at a number. Longer thinking
changes the count of orders (Listing~\ref{alg:prereq-build-orders-think}).

\pdrunlines[caption={Every order when a worker thinks for 8
turns.},label={alg:prereq-build-orders-think}]{orders-think-8}{1-8}{1\textendash 8}

A round that overlaps loses an update with chance 255/256 under the blind
design. Looking twice cuts that to 12155/65536, a little under 1/5.
\emph{Now you try:} Exercise~\ref{ex:prereq-build-walk} derives both
chances and shows how slowly the second one falls as the workers think
longer. With 1 round in 50 overlapping, a round loses an update with chance
about 0.020 for the blind design and 0.0037 for looking twice. The checked
design has no order that loses one.

\pdkernelname{} runs 1,000 such jobs for each design. Every job draws from
its own stream of a seeded generator, so the command prints the same table
each time it is run, and Figure~\ref{fig:prereq-survival} draws those
tables as they were printed.

% Reader question: if each round has only a small chance of losing an update,
% how many long jobs finish with none lost? And does a worker that checks
% its own read change the answer?
% Claim: with any fixed chance per round, the share of safe jobs falls
% geometrically to zero. Looking twice slows the fall and does not stop it.
% A check at the write keeps every job safe.
% Evidence: 1,000 seeded jobs of 600 rounds per design, counted every 25
% rounds, with the exact curve from the workload and every order of turns
% [internal, harbor survival; sessions survival-*.txt, plotted as printed].
% Must distinguish: counted from exact; the three designs; in the right
% panel, the theorem's bound from the exact curve under it, and the bursty
% curve from the steady one, which keeps its grey from the left panel.
% Grammar chosen: two panels on one vertical scale; a line for what is
% computed, dots with interval bars for what was counted; each line labelled
% at the line.
% Grammar rejected: a logarithmic scale, which makes the curves straight and
% cannot show a share of zero or a design that stays at one.
% Five-second test: which line never leaves the top?
\begin{figure}[H]
\centering
\pdfigurehue{pdslate}%
\begin{tikzpicture}[pd figure]
  \pgfplotsset{
    counted/.style={only marks,mark size=1.15pt,
      error bars/y dir=both,error bars/y explicit,
      error bars/error bar style={line width=.5pt},
      error bars/error mark options={rotate=90,mark size=1.1pt,line width=.5pt}},
  }
  % The tables are the recorded runs, set as they were printed: the first
  % line is the command, and every remark starts with '#'.
  \pgfplotstableset{
    as printed/.style={skip first n=1,x=rounds},
    bars/.style={y=share,y error plus expr=\thisrow{high}-\thisrow{share},
      y error minus expr=\thisrow{share}-\thisrow{low}},
  }
  \begin{axis}[pd axis,at={(0,0)},width=3.85cm,height=4.6cm,
    xmin=0,xmax=600,ymin=0,ymax=1.04,
    xtick={0,200,400,600},ytick={0,0.25,0.5,0.75,1},yticklabels={0,.25,.5,.75,1},
    title={steady overlap},xlabel={rounds},ylabel={share of jobs with no update lost}]
    \addplot[pd muted series,smooth] table[as printed,y=exact] {excerpts/survival-blind.txt};
    \addplot[counted,draw=pdinkmuted!70,mark options={fill=pdinkmuted!70,draw=pdpage,line width=.3pt}]
      table[as printed,bars] {excerpts/survival-blind.txt};
    \addplot[draw=pdink,smooth,mark=none] table[as printed,y=exact] {excerpts/survival-look-twice.txt};
    \addplot[counted,draw=pdink,mark options={fill=pdink,draw=pdpage,line width=.3pt}]
      table[as printed,bars] {excerpts/survival-look-twice.txt};
    \addplot[draw=pdfocus,line width=1.6pt,mark=none] table[as printed,y=exact] {excerpts/survival-checked.txt};
    \addplot[counted,draw=pdfocus,mark options={fill=pdfocus,draw=pdpage,line width=.3pt}]
      table[as printed,bars] {excerpts/survival-checked.txt};
    \node[pd focus label,anchor=north east] at (axis cs:600,0.955) {checked};
    \node[pd label,anchor=south west] at (axis cs:255,0.40) {look twice};
    \node[pd label,text=pdinkmuted,anchor=south west] at (axis cs:75,0.23) {blind};
  \end{axis}
  \begin{axis}[pd axis,at={(5.1cm,0)},width=3.85cm,height=4.6cm,
    xmin=0,xmax=600,ymin=0,ymax=1.04,
    xtick={0,200,400,600},ytick={0,0.25,0.5,0.75,1},yticklabels={},
    title={the same overlap, in bursts},xlabel={rounds}]
    \addplot[pd guide,smooth,mark=none] table[as printed,y=bound] {excerpts/survival-bursty.txt};
    \addplot[pd muted series,smooth] table[as printed,y=exact] {excerpts/survival-blind.txt};
    \addplot[draw=pdink,smooth,mark=none] table[as printed,y=exact] {excerpts/survival-bursty.txt};
    \addplot[counted,draw=pdink,mark options={fill=pdink,draw=pdpage,line width=.3pt}]
      table[as printed,bars] {excerpts/survival-bursty.txt};
    \node[pd label,anchor=south west] at (axis cs:300,0.56) {the bound};
    \node[pd label,anchor=south west] at (axis cs:118,0.25) {blind, in bursts};
  \end{axis}
\end{tikzpicture}
\caption{\textbf{A small chance of loss in every round leaves almost no long
job safe, and a chance of zero leaves every job safe.} The share of 1,000
simulated jobs with no update lost yet. Lines are exact; dots are counts with 95\%
intervals. On the right, the grey line repeats the blind curve from the
left, and the line above both curves is the bound of
Theorem~\ref{thm:non-determinism}. [internal, \texttt{harbor survival}]}
\label{fig:prereq-survival}
\end{figure}


Read the left panel first. Under the blind design the chance that a job is
still safe is a half or less from round 35, and no job of the 1,000 is safe
at round 600. Looking twice moves the halfway point to round 187 and leaves
116 jobs safe at round 600. It bought about five times as long and the same
end. Under the checked design all 1,000 jobs are safe at every row.
Listing~\ref{alg:prereq-build-survival} is part of the table behind the
blind curve.

\pdrunlines[caption={Part of what \texttt{harbor survival} prints for the
blind design. The figure plots this
file.},label={alg:prereq-build-survival}]{survival-blind}{1,6-7,14-19,39-40}{1, 6\textendash 7, 14\textendash 19 and 39\textendash 40}

\paragraph{Why simulate what can be worked out.}
The lines in the figure were not fitted to the dots. The chance of loss in
one round is exact, so the share of jobs safe after $h$ rounds is
$(1-\epsilon)^h$, with equality because here every round has the same chance
whatever came before. The simulation was run for another reason. It
executes the workers of Listings~\ref{alg:prereq-build-blind}
to~\ref{alg:prereq-build-checked} and the guard of
Listing~\ref{alg:prereq-build-current}, turn by turn, and so it would expose a
mistake in that code, which the arithmetic could never see. Each checks the
other, and they agree. The exact value lies inside the counted interval in
23 of the 25 rows for the blind design and in all 25 for looking twice. A
95\% interval is built to miss about one row in twenty, so one or two misses
in 25 rows are what agreement looks like.

\paragraph{What the dots cannot show.}
The bars are Wilson's 95\% intervals~\pdcite{wilson1927probable}. The usual
interval, the share plus or minus two standard errors, shrinks to a single
point when every job comes out the same way, and it is unreliable well
before that~\pdcite{brown2001interval}. Every job coming out the same way is
exactly the case of the checked design. Wilson's interval for 1,000 safe
jobs of 1,000 runs from 0.9962 to 1. That is all the simulation supports:
fewer than four jobs in a thousand lose an update. A thousand runs cannot
show that the chance is zero, and a million could not either. The count of
every order shows it for two workers that think for 8 turns, and part~(1)
of Theorem~\ref{thm:kernel-gap} shows it for any number of workers. A sample
can bound a chance. It takes the count or the proof to make the chance zero.

\paragraph{When the rounds are not alike.}
The theorem does not ask for independent steps, and the right panel is a
workload that has none. The overlap is still 1 round in 50 on average, and
it arrives in bursts. In a calm spell 1 round of 500 overlaps, and in a busy
spell 1 round of 5. A calm spell lasts 50 rounds on average and a busy one
5, so whether a round overlaps depends on the rounds before it.

The theorem's premise is a floor: after any safe history, the chance of loss
in the next step is at least $\epsilon$. The floor here is the chance in a
calm round, and the dotted line is the bound it gives. The blind design
falls far below it. At round 600 the bound allows a share of 0.30, and the
exact share is 0.0005. A bound that knows only the calmest rounds is loose
when busy rounds are common.

The grey line is the steady workload again, with the same average overlap,
and it lies below the bursty curve. At round 100 the exact shares are 0.13
when the overlap is steady and 0.29 when it comes in bursts. A job can be
lost only once, so overlaps crowded into a burst are partly spent on jobs
that are lost already. Someone who puts the average chance into the
theorem's product gets the grey line, and the truth is above it. An average
is not a floor.
\pdscope{The curves describe one model of a workload. Change the scheduler
or the thinking time and every number in the figure moves. Two conclusions
do not move: a chance of loss above zero in every round empties the jobs,
and a chance of zero keeps them.}

\begin{pdboundary}
\setlength{\parskip}{4pt plus 1pt}
The cell is one number held in memory. Nothing here survives a crash, so
part~(3) of Theorem~\ref{thm:kernel-gap} is not built in this section.
Chapter~1 keeps its state in a database and kills the process to test it.
The checked design promises that no update is lost. It does not promise
every worker a turn. A worker can be refused again and again while other
writes land, and Herlihy's name for a design that rules this out is
wait-free~\pdcite{herlihy1991waitfree}. The checked cell is not wait-free:
each refusal means that some write landed, so the workers as a group always
move forward, and one of them can still starve. A version over one number is
also easy to define. Real work reads several things and changes one, and
Exercise~\ref{ex:prereq-build-skew} shows what a version on each thing
separately lets through.
\end{pdboundary}
